\chapter{Hardware and Platform Overview}

This chapter provides an overview of the hardware and cloud platform used in our project. We introduce the Yolo Uno development board and the Core IOT platform, outlining the setup process and dashboard design. Detailed firmware implementation and code explanations are reserved for Chapter 7.

\section{Introduction to Yolo Uno}

The Yolo Uno is a versatile development board designed for IoT applications, featuring the powerful ESP32 microcontroller. It integrates essential components such as Wi-Fi, Bluetooth, and various onboard sensors/actuators, making it an ideal choice for rapid prototyping.

Key features of the Yolo Uno include:
\begin{itemize}
    \item \textbf{ESP32-WROOM-32E}: Dual-core processor with Wi-Fi and Bluetooth connectivity.
    \item \textbf{Onboard Peripherals}: Includes NeoPixel LEDs, buttons, and Grove connectors for easy sensor expansion.
    \item \textbf{USB-C Interface}: For programming and power.
    \item \textbf{Compact Form Factor}: Suitable for embedded projects.
\end{itemize}

\begin{figure}[H]
    \centering
    \includegraphics[width=0.6\textwidth]{assets/yolouno.png}
    \caption{The Yolo Uno development board.}
    \label{fig:yolouno}
\end{figure}

\section{Core IOT Platform Setup}

Core IOT is a cloud-based IoT platform built on the ThingsBoard framework. We selected this platform for our project due to its open-source nature and comprehensive feature set for device management and data visualization.

\subsection{Registration Process}

We accessed the Core IOT platform at \url{https://coreiot.io} and clicked the "Start Now" button to begin registration (Figure~\ref{fig:coreiot_homepage}).

\begin{figure}[H]
    \centering
    \includegraphics[width=1\textwidth]{assets/Coreiot.png}
    \caption{Core IOT homepage with the "Start Now" button for registration.}
    \label{fig:coreiot_homepage}
\end{figure}

The platform offers multiple login options through OAuth providers. We chose to authenticate using a Google account for convenience (Figure~\ref{fig:coreiot_login}).

\begin{figure}[H]
    \centering
    \includegraphics[width=1\textwidth]{assets/Login.png}
    \caption{Login interface with OAuth provider options.}
    \label{fig:coreiot_login}
\end{figure}

After successful authentication, we were assigned the Tenant Administrator role, which provides full access to device management and dashboard creation features (Figure~\ref{fig:coreiot_dashboard}).

\begin{figure}[H]
    \centering
    \includegraphics[width=1\textwidth]{assets/platform.png}
    \caption{Dashboard view after login showing Tenant Administrator role.}
    \label{fig:coreiot_dashboard}
\end{figure}

\subsection{Device Provisioning}

With the account ready, we proceeded to register our ESP32-based device on the platform. In the "Devices" section, we created a new device and obtained an access token. This token is required for the device to authenticate when sending telemetry data via MQTT.

To verify the connection and data transmission capabilities, we utilized a Python script to simulate the IoT device behavior before deploying it to the actual hardware. It is critical to note that we specifically selected the \texttt{paho-mqtt} library version 1.6.1 for this implementation, as newer versions (e.g., 2.0) introduce breaking changes in the authentication handling that are incompatible with our current configuration.

\section{Telemetry Data Transmission}

After configuring the device credentials, we developed a firmware routine to format sensor data into a JSON structure compliant with the Core IOT platform's requirements. The main task runs in a FreeRTOS task, waiting for WiFi connectivity before attempting to connect to Core IOT. It then enters a loop that publishes telemetry data every 10 seconds.

We verified the data ingestion by navigating to the "Latest Telemetry" tab within the device details page on the Core IOT portal. This interface allows us to confirm that the JSON payload is being parsed correctly and that values for temperature and humidity are updating dynamically as sent by the device.

\begin{figure}[H]
    \centering
    \includegraphics[width=0.9\textwidth]{assets/lastest_telemetry.png}
    \caption{The "Latest Telemetry" tab displaying real-time data received from the device.}
    \label{fig:latest_telemetry}
\end{figure}

\section{Dashboard Implementation}

One of the primary advantages of the Core IOT platform is its capability to visualize data through customizable dashboards. We created a comprehensive monitoring dashboard named \textit{Control and Monitoring center} to serve as the central interface for real-time system oversight and actuator control. The dashboard employs a multi-panel layout that organizes information hierarchically, prioritizing critical control functions and environmental data visualization.

\subsection{Dashboard Layout Architecture}

We designed the dashboard with a responsive grid layout divided into three primary horizontal sections. The top section contains system status indicators and actuator controls, providing immediate access to the most frequently used functions. The middle section displays real-time sensor data through time-series charts and gauge widgets, enabling trend analysis and anomaly detection. The bottom section presents an alarm management table and geographical visualization, supporting incident tracking and multi-device deployment scenarios.

This hierarchical organization follows established human-computer interaction principles, placing high-priority interactive elements within the primary visual field while maintaining comprehensive data visibility through strategic use of white space and visual grouping.

\subsection{Status Indicators and Device Metrics}

The dashboard header displays two key performance indicators using card widgets. The first card shows the active device count, labeled \textbf{Devices}, which indicates the number of ESP32 devices currently connected to the platform. In our deployment, this displays a value of 1, confirming that our primary environmental monitoring device maintains an active MQTT connection.

Adjacent to the device count, we positioned a \textbf{Total} card that aggregates a custom metric across all connected devices. This counter widget can be configured to track various cumulative statistics such as total alarm count, cumulative uptime, or aggregated sensor readings. The card widgets utilize ThingsBoard's entity aggregation functions, which query the time-series database to compute values based on configurable time windows.

To the right of the status cards, we integrated a WiFi quality indicator using an icon widget. This widget dynamically changes color based on the RSSI (Received Signal Strength Indicator) value transmitted in the device telemetry. Strong signal strength displays a green icon, while degraded connectivity transitions to yellow or red, providing immediate visual feedback on network health without requiring detailed inspection of raw telemetry values.

\subsection{Remote Control Panel}

We implemented a dedicated control panel section containing three RPC toggle switches for actuator management: Water Pump, Fan, and NeoLed (RGB LED strip). Each switch widget is configured with bidirectional RPC methods to both read the current state and issue control commands.

The switch configuration employs getter and setter RPC methods. When the dashboard loads, each switch automatically invokes its corresponding getter method (\texttt{getFanStatus}, \texttt{getWaterPumpStatus}, \texttt{getLedStatus}) to synchronize the UI state with the actual device state. This prevents desynchronization scenarios where the dashboard might display \textbf{OFF} while the device is actively driving the actuator.

When a user toggles a switch, the dashboard transmits an RPC request containing the method name and the desired boolean state. For example, clicking the Fan switch from OFF to ON triggers an RPC call to \texttt{setFanStatus} with parameter \texttt{true}. The device firmware processes this command, actuates the corresponding GPIO pin, and returns an \texttt{RPC\_Response} to acknowledge successful execution.

We observed that timeout errors occasionally occur when the device is processing intensive tasks or when network latency exceeds the default RPC timeout threshold (5 seconds). In such cases, the switch widget displays a \textit{Request timeout} modal, prompting the user to retry the command. To mitigate this issue, we optimized the RPC callback functions to execute GPIO operations quickly and defer any blocking operations to separate FreeRTOS tasks.

\subsection{Environmental Data Visualization}

The central focus of the dashboard is the environmental monitoring section, which utilizes time-series charts to visualize temperature and humidity trends. We configured a line chart widget with dual Y-axes, where temperature is plotted in degrees Celsius on the left axis and relative humidity percentage on the right axis. This dual-axis configuration enables direct visual correlation between the two variables, which is essential for identifying environmental patterns such as evaporative cooling effects.

The chart widget subscribes to the \texttt{temperature} and \texttt{humidity} telemetry keys and maintains a rolling time window of the last hour by default. Users can adjust this window through the chart's time range selector, enabling both real-time monitoring and historical analysis. We customized the chart appearance with distinct colors: red for temperature and blue for humidity, following conventional meteorological visualization standards.

Below the time-series chart, we positioned a second chart dedicated to soil moisture monitoring. This widget displays the \texttt{moisture} telemetry key, providing agricultural and plant care applications with a clear view of substrate hydration levels over time. The moisture chart uses a green color scheme and includes threshold lines at 30\% (dry) and 70\% (saturated) to guide irrigation decisions.

Adjacent to the charts, we implemented a large numeric display widget showing the current light level in lux. This widget is configured with the \texttt{light} telemetry key and formatted to display values with scientific notation (e.g., 65000 m², though the unit should technically be lux). The large font size makes this value readable at a glance, which is particularly useful for automated greenhouse or indoor farming applications where lighting optimization is critical.


\subsection{Widget Configuration Best Practices}

Throughout the dashboard design process, we identified several critical configuration requirements:

\textbf{Data Key Consistency:} The telemetry keys defined in the firmware (\texttt{temperature}, \texttt{humidity}, \texttt{light}, etc.) must exactly match the data keys configured in each widget. Case sensitivity and spelling errors will result in widgets displaying \textit{No data} even when the device is actively transmitting.

\textbf{RPC Method Naming:} Switch widgets require precise RPC method names that correspond to the registered callbacks in the firmware. We standardized on a naming convention using camelCase with \textit{set} and \textit{get} prefixes (e.g., \texttt{setFanStatus}, \texttt{getFanStatus}).

\textbf{Widget Refresh Rates:} We configured chart widgets to poll for new data every 5 seconds, balancing real-time responsiveness with server load. Faster refresh rates increase database query frequency, which can impact platform performance in deployments with hundreds of devices.

\textbf{Color Coding and Visual Hierarchy:} We applied consistent color schemes throughout the dashboard: green for normal conditions, yellow for warnings, and red for critical alerts. This visual language enables rapid situation assessment without requiring detailed reading of all numeric values.

\begin{figure}[H]
    \centering
    \includegraphics[width=1\textwidth]{assets/control_monitoring_dashboard.png}
    \caption{Complete Control and Monitoring center dashboard showing status cards, RPC switches, time-series charts, light level display, WiFi indicator, alarm table, and map widget.}
    \label{fig:final_dashboard}
\end{figure}

\section{Summary}

In this chapter, we introduced the Yolo Uno hardware and the Core IOT cloud platform. We outlined the initial setup, including account registration, device provisioning, and the conceptual design of the monitoring dashboard. The specific software implementation, including MQTT communication and FreeRTOS tasks, will be detailed in Chapter 7.
